\section{실험 결과}

\subsection*{실패 발견율과 Random 테스트 대비 향상}
\begin{table*}[!t]
\centering\small
\caption{AFS와 Random 테스트의 실패 발견 결과}\label{tab:result}
\begin{tabularx}{\textwidth}{@{}Ycccc@{}}
\toprule
& \multicolumn{2}{c}{실험 1: 로봇팔} & \multicolumn{2}{c}{실험 2: 휴머노이드} \\
& \multicolumn{2}{c}{(3세트 $\times$ 2회)} & \multicolumn{2}{c}{(3세트 $\times$ 6회)} \\
\cmidrule(lr){2-3}\cmidrule(lr){4-5}
& AFS & Random 테스트 & AFS & Random 테스트 \\
\midrule
시도 / 유효 / 제외 & 6 / 6 / 0 & 6 / 6 / 0 & 18 / 18 / 0 & 18 / 18 / 0 \\
PASS / FAIL & 0 / 6 & 5 / 1 & 12 / 6 & 16 / 2 \\
\textbf{실패 발견율 (FDR)} & \textbf{100\%} & 16.7\% & \textbf{33.3\%} & 11.1\% \\
실패가 나온 서로 다른 환경 수 & 3 & 1 & 6 & 2 \\
\midrule
사람 안전 위험 검출 (그중 작업 성공) & -- & -- & 11/18 (9) & 14/18 (12) \\
AFS 제안 실행 수 & 6 & -- & 18 & -- \\
총 실행 시간 / 로봇 판단 횟수 & \multicolumn{2}{c}{약 90초 / 40초} & \multicolumn{2}{c}{판단 196회 / 147회} \\
\bottomrule
\end{tabularx}
\end{table*}

\begin{figure}[!t]
\centering
\begin{tikzpicture}
\begin{axis}[
  ybar, width=\linewidth, height=48mm, bar width=9pt,
  ymin=0, ymax=125, ylabel={FDR (\%)}, ylabel style={font=\small},
  symbolic x coords={A,B}, xtick=data,
  xticklabels={{로봇팔},{휴머노이드}},
  xticklabel style={font=\scriptsize,align=center},
  tick label style={font=\scriptsize},
  nodes near coords, nodes near coords style={font=\scriptsize},
  legend style={font=\scriptsize,at={(0.5,1)},anchor=north,draw=none,legend columns=2},
  axis lines*=left, enlarge x limits=0.4]
\addplot[fill=afs,draw=afs] coordinates {(A,100) (B,33.3)};
\addplot[fill=rnd,draw=rnd] coordinates {(A,16.7) (B,11.1)};
\legend{AFS, Random 테스트}
\end{axis}
\end{tikzpicture}
\caption{AFS와 Random 테스트의 실패 발견율. 두 실험 모두 각 세트의 공통 시작 장면을 제외한 평가 실행 기준이다.}\label{fig:fdr}
\end{figure}

실험 결과를 표~\ref{tab:result}와 그림~\ref{fig:fdr}에 정리하였다. 
로봇팔 실험은 3세트에서 각 방법을 2회씩, 방법당 6회 평가하였다. 
각 세트의 처음 2회는 평가 전에 두 방법이 같은 장면을 쓰는 공통 시작 장면으로 평가에 포함하지 않는다. 
AFS는 6회 모두 실패를 찾아 FDR 100\%로 목표 30\%를 넘었고, Random 테스트(16.7\%, 1/6) 대비 상대 향상률은 $(100-16.7)/16.7 = 500\%$로 목표 20\%를 크게 웃돈다. 
세트별로도 AFS는 2/2/2건, Random 테스트는 0/1/0건으로 모든 세트에서 AFS가 더 많은 실패를 발견했다.

휴머노이드 실험은 3세트에서 각 방법을 6회씩, 방법당 18회 평가하였다(제외 0회). 
각 세트에서는 평가 전에 두 방법이 공유하는 공통 시작 장면을 1회 실행하였으며(모두 PASS), 방법 간 차이를 만들지 않으므로 평가에 포함하지 않는다.
AFS의 FDR은 33.3\%(6/18), Random 테스트는 11.1\%(2/18)로 AFS가 3배 많은 실패 사례를 찾았고 상대 향상률은 $(33.3-11.1)/11.1 = 200\%$다.
세트별로 AFS는 세 세트 모두 2건씩(표준편차 0\%p) 실패를 찾은 반면, Random 테스트는 0/0/2건으로 세트에 따라 편차가 컸다. 
예비 실험(1세트, 방법당 유효 4회)에서도 AFS 50\%(2/4), Random 테스트 0\%(0/4)로 같은 경향이었다. 
사람 안전 위험은 AFS 11/18회, Random 테스트 14/18회에서 검출되었다.
실패한 실행은 제한 시간까지 이어지므로, 두 실험 모두 AFS의 총 실행 시간과 로봇 판단 횟수가 Random 테스트 보다 많다.

LIBERO는 Spatial·Object·Goal·Long 등 여러 과제 묶음으로 구성되고 과제마다 장면 변수를 따로 정의해야 한다.
본 과제는 Spatial의 대표 과제 하나를 주 실험으로 삼았다. 
다른 묶음에서도 같은 경향인지 확인하기 위해 GR00T~N1.7의 묶음별 공개 체크포인트로 Object(버터를 바구니에 넣기)와 Goal(서랍 열기)의 과제 1개씩을 골라 3세트$\times$2회 평가(공통 시작 장면 제외)를 독립적으로 3번 반복하였다.
Object에서 AFS는 16/18회(88.9\%), Random 테스트는 6/18회(33.3\%) 실패를 찾았고 세 번 모두 AFS가 많았다. 
다만 AFS 제안은 대부분 버터 질량과 장애물 크기를 범위 끝값으로 조합한 것이었다. 
Goal은 정책이 기본 장면에서도 거의 항상 실패하여(Random 테스트 18/18) 비교할 수 없었다.
휴머노이드도 통로 장면 외에 창고·교차로 배치에서 공통 시작 장면 외에 각 방법을 2회씩 실행하여 창고 AFS 2/2·Random 테스트 0/2, 교차로 1/2·1/2회 실패를 얻었다.
창고에서는 LLM 응답을 기다리며 멈춘 로봇이 횡단하는 사람과 접촉한 뒤 넘어졌다. 사무실 배치는 시작 장면부터 실패하여 비교하지 않았다.

\begin{table*}[!t]
\centering\scriptsize
\caption{로봇팔 실험 시도별 결과 (공통: 공통 시작 장면, 탐색: 탐색 단계, 시간: 실행 시간, 접촉: 로봇–장애물 접촉 시간)}\label{tab:trials}
\begin{tabular}{@{}lccccc|ccccc@{}}
\toprule
& \multicolumn{5}{c|}{AFS} & \multicolumn{5}{c}{Random 테스트} \\
세트(시드) / 순번 & 단계 & 결과 & 시간(s) & 접촉(s) & 그릇–접시(m) & 단계 & 결과 & 시간(s) & 접촉(s) & 그릇–접시(m) \\
\midrule
661 / 0 & 공통 & FAIL & 15.0 & 13.33 & 0.284 & 공통 & FAIL & 15.0 & 13.33 & 0.284 \\
661 / 1 & 공통 & PASS & 4.8 & 0.00 & 0.021 & 공통 & PASS & 4.8 & 0.00 & 0.021 \\
661 / 2 & 탐색 & FAIL & 15.0 & 11.04 & 0.284 & 탐색 & PASS & 4.8 & 0.00 & 0.007 \\
661 / 3 & 탐색 & FAIL & 15.0 & 10.73 & 0.284 & 탐색 & PASS & 5.1 & 0.00 & 0.011 \\
673 / 0 & 공통 & PASS & 5.2 & 0.02 & 0.019 & 공통 & PASS & 5.2 & 0.02 & 0.019 \\
673 / 1 & 공통 & PASS & 4.9 & 0.00 & 0.008 & 공통 & PASS & 4.9 & 0.00 & 0.008 \\
673 / 2 & 탐색 & FAIL & 15.0 & 12.79 & 0.275 & 탐색 & FAIL & 15.0 & 13.24 & 0.254 \\
673 / 3 & 탐색 & FAIL & 15.0 & 12.45 & 0.275 & 탐색 & PASS & 5.1 & 0.00 & 0.013 \\
677 / 0 & 공통 & PASS & 4.8 & 0.23 & 0.019 & 공통 & PASS & 4.8 & 0.23 & 0.019 \\
677 / 1 & 공통 & PASS & 5.0 & 0.13 & 0.004 & 공통 & PASS & 5.0 & 0.13 & 0.004 \\
677 / 2 & 탐색 & FAIL & 15.0 & 13.73 & 0.312 & 탐색 & PASS & 5.0 & 0.00 & 0.014 \\
677 / 3 & 탐색 & FAIL & 15.0 & 13.74 & 0.312 & 탐색 & PASS & 5.0 & 0.00 & 0.009 \\
\bottomrule
\end{tabular}
\end{table*}

표~\ref{tab:trials}는 로봇팔 실험의 시도별 기록이다. 
공통 시작 장면에서는 두 방법의 결과가 같고, 차이는 모두 탐색 단계에서 나온다. 
AFS의 탐색 단계 FAIL 6건 중 3건은 같은 세트에서 이미 실패한 장면을 다시 제안한 경우로, 새로 찾은 실패 환경은 3건이다(Random 테스트 1건).

\subsection*{AFS 예상의 적중}
AFS는 제안마다 예상 실패를 함께 내므로, 예상이 실제 결과로 이어졌는지 확인할 수 있다. 
로봇팔에서 6건의 제안은 모두 ``이송 경로의 장애물이 접근을 막을 것''이라는 가설이었고, 6건 모두 로봇–장애물 접촉(10.7--13.7초)을 동반한 FAIL로 판정되어 가설과 관측이 일치했다. 
휴머노이드 본 실험에서는 18건의 제안 중 6건이 FAIL로 이어졌다. 
목표 위치 점유를 예상한 3건은 모두 실제로 목표 위치 점유로 판정되었고, 나머지 FAIL 3건은 물체 조작 실패 2건, 유형 미부여 1건이었다. 
반면 충돌(5건)과 이동 정체(4건)를 노린 장애물 배치·바닥 마찰 변경 9건은 모두 PASS였다. 
두 방법의 FAIL 8건은 모두 상자의 위치나 통로 폭을 바꾼 경우였으므로, 이 로봇 제어 구성의 취약점은 장애물 형상이나 마찰보다 \textbf{목표 부근 상자 배치와 통로 폭}에 있다.
AFS 가설은 관측된 실패 메커니즘(점유, 차단)을 겨냥할 때 잘 맞았고, 아직 관측되지 않은 유형을 겨냥할 때는 빗나갔다.

\subsection*{실패 유형 다양성}
\begin{table*}[!t]
\centering\small
\caption{초기 정의 6종 기준 실패 유형 발견 현황}\label{tab:div}
\begin{tabularx}{\textwidth}{@{}p{24mm}cp{26mm}Y@{}}
\toprule
유형 & 발견 & 실험 & 관측 근거 \\
\midrule
충돌 & \checkmark & 로봇팔 & 로봇–장애물 접촉 10.7--13.7초, 최대 접촉력 약 230--380\,N, 그릇 이동 1\,mm 미만 \\
장애물 간섭 & \checkmark & 휴머노이드 & 장애물 부근에서 경로 계획이 반복 실패하며 정체 \\
목표 위치 점유 & \checkmark & 휴머노이드 & 상자가 목표 영역을 점유, 목표 0.86\,m 앞에서 정체 \\
작업자 안전 위험 & \checkmark & 휴머노이드 & 사람과 1\,m 이내 0.5초 이상 노출 또는 접촉, AFS 11/18회 \\
도달 불가 & 부분 & 두 실험 & 그릇을 들어 운반했으나 접시 8.7\,cm 앞에서 미완료, 최종 목표점 경로 없음 반복 \\
인식 오류 & -- & -- & 가림·유사 물체 조건은 후속 실험 \\
\midrule
\multicolumn{4}{@{}l}{\textbf{합계: 6종 중 4종 발견 + 도달 불가 부분 관측 (목표 4종 달성)}} \\
\bottomrule
\end{tabularx}
\end{table*}
표~\ref{tab:div}는 초기 정의 6종 기준의 발견 현황이다. 
작업자 안전 위험은 목표 달성 여부와 무관하게 안전 요구사항을 위반한 실행을 실패 사례로 본다.
실제 실행에서는 초기 6종으로 설명하기 어려운 행동 실패도 관측되어, 넘어짐·이동 정체·물체 조작 실패·목표 완료 실패와 시스템 연동 문제 2종(목표점 경로 없음, 중간 목표 반복 도착)을 더한 12종 체계로 확장하였다(표~\ref{tab:taxonomy}).
이 확장 체계로는 목표 완료 실패, 물체 조작 실패, 중간 목표 반복과 창고 장면의 넘어짐을 추가로 관측하였다.
